\documentclass[12pt,a4paper]{article}
\usepackage[UTF8]{ctex}
\usepackage{geometry}
\usepackage{enumitem}
\usepackage{listings}
\usepackage{xcolor}
\usepackage{hyperref}
\usepackage{titlesec}
\usepackage{fontspec}
\usepackage{longtable}

\geometry{left=2.5cm, right=2.5cm, top=2.5cm, bottom=2.5cm}

% 代码样式
\lstset{
    basicstyle=\small\ttfamily,
    backgroundcolor=\color{gray!10},
    frame=single,
    rulecolor=\color{gray!30},
    breaklines=true,
    breakatwhitespace=false,
    columns=flexible,
    keepspaces=true,
    tabsize=4,
    showstringspaces=false,
    numberstyle=\tiny\color{gray},
    keywordstyle=\color{blue!70},
    commentstyle=\color{green!50!black},
    stringstyle=\color{red!60!black},
    numbers=left,
    numbersep=5pt,
    xleftmargin=2em,
    framexleftmargin=1.5em
}

\setlength{\parskip}{4pt}
\setlength{\parindent}{2em}

\title{\huge \textbf{Git 学习笔记}}
\author{}
\date{\today}

\begin{document}

\maketitle
\thispagestyle{empty}
\newpage

\tableofcontents
\newpage

%====================================================================
\section{引言 —— 为什么要用 Git？}
%====================================================================

\subsection{什么是版本控制？}

假设你在写毕业论文，你一定经历过这样的场景：

\begin{lstlisting}[frame=none, numbers=none]
毕业论文_初稿.docx
毕业论文_修改1.docx
毕业论文_修改2.docx
毕业论文_最终版.docx
毕业论文_真的最终版.docx
毕业论文_打死也不改了.docx
\end{lstlisting}

这就是\textbf{没有版本控制}的混乱状态。版本控制系统（VCS）就是为了解决这个问题而生——
它帮你记录每次修改的历史，让你随时可以回到任意版本，还能多人并行协作而不互相覆盖。

\subsection{集中式 vs 分布式}

\begin{longtable}{|p{4cm}|p{5cm}|p{5cm}|}
\hline
\textbf{特性} & \textbf{集中式（SVN、CVS）} & \textbf{分布式（Git）} \\
\hline
仓库位置 & 只有一个中央服务器 & 每个人本地都有完整仓库 \\
\hline
离线工作 & 不行，必须联网 & 可以，大部分操作本地完成 \\
\hline
速度 & 慢（大部分操作需要网络） & 快（本地操作） \\
\hline
安全性 & 中央服务器挂了就全丢了 & 每个人都是备份 \\
\hline
分支代价 & 高（完整复制整个项目） & 极低（只是一个指针） \\
\hline
\end{longtable}

\subsection{Git 的诞生}

2005 年，Linux 内核社区和当时的版本控制系统 BitKeeper 闹翻了。Linus Torvalds 花了两周时间，用 C 语言写出了 Git。设计目标是：
\begin{itemize}[itemsep=3pt]
    \item 速度飞快
    \item 设计简单（数据结构精妙但简洁）
    \item 完全分布式
    \item 支持非线性开发（成千上万个分支）
\end{itemize}

如今，Git 已经是全世界最流行的版本控制系统，无论个人项目还是企业级开发，几乎都在用。

%====================================================================
\section{基本概念 —— 理解 Git 的三个区域}
%====================================================================

\subsection{工作区、暂存区与仓库}

这是 Git 最核心的抽象，理解了这三个区域，你就理解了 Git 的一半。

\begin{itemize}[itemsep=5pt]
    \item \textbf{工作区（Working Directory）}
    \begin{itemize}
        \item 就是你电脑上看到的文件目录。
        \item 你在这里编辑文件、写代码、改 bug。
        \item 这是你可以直接看到、操作的区域。
    \end{itemize}

    \item \textbf{暂存区（Staging Area / Index）}
    \begin{itemize}
        \item 可以理解为一个\textbf{购物车}。
        \item 你改了一堆文件，但不想一次性全提交。
        \item 用 \texttt{git add} 把想要提交的文件先放进购物车。
        \item 暂存区让你可以\textbf{精心挑选}每次提交的内容。
    \end{itemize}

    \item \textbf{本地仓库（Repository / .git）}
    \begin{itemize}
        \item 就是隐藏的 \texttt{.git/} 目录。
        \item 存着所有的版本历史、分支信息、标签等。
        \item 相当于 Git 的\textbf{数据库}。
        \item \texttt{git commit} 就是把暂存区的内容写入这个数据库。
    \end{itemize}
\end{itemize}

\subsubsection{完整工作流程}

\begin{lstlisting}[numbers=none]
初始状态                           git add 之后              git commit 之后
工作区: 修改了 A B C 三个文件  -->  暂存区: A B C           暂存区: (空)
                                   工作区: A B C           本地仓库: (记录了一次新版本)
                                                           工作区: A B C
\end{lstlisting}

\textbf{打个比方}：你在做饭（工作区），切好的菜放进盘子里（暂存区\ --- 随时可以再加或拿走），
准备好后一起下锅炒（提交）。每炒好一道菜就是一个版本。

\subsection{Git 对象模型}

Git 本质上是一个\textbf{内容寻址的文件系统}，它用哈希值（SHA-1）来定位所有内容。

\subsubsection{四种核心对象}

\begin{enumerate}[itemsep=5pt]
    \item \textbf{blob（Binary Large Object）}
    \begin{itemize}
        \item 存储\textbf{文件内容}。
        \item 不包含文件名，只是内容的二进制快照。
        \item 两个文件即使名字不同，如果内容相同，指向同一个 blob。
    \end{itemize}

    \item \textbf{tree（树对象）}
    \begin{itemize}
        \item 存储\textbf{目录结构}。
        \item 记录文件名 $\rightarrow$ blob 的映射，以及子目录 $\rightarrow$ tree 的映射。
        \item 相当于文件系统中的目录（文件夹）。
    \end{itemize}

    \item \textbf{commit（提交对象）}
    \begin{itemize}
        \item 存储一次\textbf{提交快照}。
        \item 包含：指向一棵 tree、父提交的哈希、作者信息、时间戳、提交消息。
        \item 每个 commit 就是一个\textbf{版本}。
    \end{itemize}

    \item \textbf{tag（标签对象）}
    \begin{itemize}
        \item 为某个 commit 起一个\textbf{固定的名字}（如 v1.0.0）。
        \item 分支会移动，但标签不会。
    \end{itemize}
\end{enumerate}

\subsubsection{对象之间的关系}

\begin{lstlisting}[numbers=none]
commit (abc123)
  |-- 指向 tree (def456)
  |-- 父提交: 上一个 commit
  |-- 作者/时间/消息
        |
        tree (def456)
        |-- "main.c"    -->  blob (aaa111)   # 文件内容
        |-- "util.h"    -->  blob (bbb222)
        |-- "src/"      -->  tree (ccc333)    # 子目录
\end{lstlisting}

所有的对象都通过 40 位的 SHA-1 哈希值标识，例如 \texttt{e86a3764c7a5b5e8a440b8e139e4c0e6c2e4d6b8}。
Git 保证：只要内容不同，哈希值就不同。所以哈希值既是地址，也是校验和。

%====================================================================
\section{基本操作 —— 从零开始用 Git}
%====================================================================

\subsection{初始化与配置}

\subsubsection{第一次使用 Git 的必做配置}

\begin{lstlisting}
# Git 需要知道你是谁，因为每次提交都会记录作者信息
git config --global user.name "Your Name"
git config --global user.email "email@example.com"

# 查看所有配置
git config --list

# 建议：设置默认分支名为 main（旧版 Git 默认是 master）
git config --global init.defaultBranch main
\end{lstlisting}

配置文件有三个层级：
\begin{itemize}[itemsep=3pt]
    \item \texttt{--system}：系统级，影响所有用户（很少用）
    \item \texttt{--global}：用户级，影响你的所有仓库（最常用）
    \item \texttt{--local}：仓库级，只影响当前仓库（默认）
\end{itemize}
优先级：local > global > system。

\subsection{初始化仓库}

\begin{lstlisting}
# 进入你的项目目录
cd my-project

# 把当前目录变成 Git 仓库（会创建一个 .git/ 隐藏目录）
git init

# 执行后，你就可以开始用 Git 管理这个目录的文件了
\end{lstlisting}

\texttt{git init} 只是在当前目录创建了一个 \texttt{.git/} 文件夹，
你的文件本身没有任何变化，Git 从此开始监视这个目录。

\subsection{文件的生命周期}

每个文件在 Git 中有四种状态：

\begin{lstlisting}[numbers=none]
Untracked（未跟踪）
  |-- 文件刚创建，Git 还没有记录它
  |-- git status 会显示为红色
  |-- git add 后变成 Staged
  |
Tracked（已跟踪）下的三种子状态：
  |-- Unmodified（未修改）：和仓库里的内容一样
  |-- Modified（已修改）：改了但还没 add
  |-- Staged（已暂存）：add 了但还没 commit

状态转换图：
Untracked  --git add-->  Staged  --git commit-->  Unmodified
Modified   --git add-->  Staged
Unmodified --修改文件-->  Modified
\end{lstlisting}

\subsection{查看状态 —— git status}

这是\textbf{最常用的命令}，没有之一。任何时候不确定当前状态，先敲它。

\begin{lstlisting}
# 假设你刚 init，目录里只有一个 README.md
$ git status
On branch main

No commits yet

Untracked files:
  (use "git add <file>..." to include in what will be committed)
        README.md

nothing added to commit but untracked files present

# 简洁模式
$ git status -s
?? README.md          # ?? 表示 Untracked
\end{lstlisting}

状态符号含义：
\begin{itemize}[itemsep=3pt]
    \item \texttt{??} — 未跟踪
    \item \texttt{M} — 已修改
    \item \texttt{A} — 新添加到暂存区
    \item \texttt{D} — 已删除
    \item \texttt{R} — 已重命名
    \item 左列 = 暂存区状态，右列 = 工作区状态
\end{itemize}

\subsection{添加文件 —— git add}

\begin{lstlisting}
# 添加单个文件
git add index.html

# 添加所有变更（新增、修改、删除）
git add .

# 添加当前目录的所有 .py 文件
git add *.py

# 交互式逐块添加（可以只提交文件的一部分修改！）
git add -p
# 这个命令会逐个"块"询问你是否要暂存
# y = 暂存这一块, n = 跳过, s = 再拆细, e = 手动编辑
\end{lstlisting}

\texttt{git add -p} 非常实用。比如你改了某个文件，有些修改是关于修复 bug A，
有些是添加功能 B，你可以用 \texttt{git add -p} 只暂存 A 的部分，先提交 A，
再暂存 B 的部分，提交 B。做到\textbf{一次提交只做一件事}。

\subsection{提交 —— git commit}

\begin{lstlisting}
# 最简单的提交
git commit -m "修复了登录页面的 bug"

# 提交前自动 add 所有已跟踪文件的修改（不包含新文件）
git commit -a -m "更新配置文件和文档"

# 修改最近一次提交（如果漏了文件或提交消息写错了）
git commit --amend
# 注意：会改变提交的哈希值，不要用于已推送的提交！
\end{lstlisting}

\textbf{提交的最佳实践}：
\begin{itemize}[itemsep=3pt]
    \item 每次提交只做一件事，保持原子性。
    \item 提交消息要说明\textbf{为什么}修改，而不是\textbf{改了什么}。
    \item 坏例子：\texttt{修改了 login.py}
    \item 好例子：\texttt{修复了密码输入过长时的截断 bug}
\end{itemize}

\subsection{完整的新项目流程示例}

\begin{lstlisting}
# 1. 初始化
git init
git config --global user.name "张三"
git config --global user.email "zhangsan@example.com"

# 2. 创建文件并提交
echo "# 我的项目" > README.md
git add README.md
git commit -m "docs: 初始化项目说明"

# 3. 写了一些代码后再提交
git add main.py utils.py
git commit -m "feat: 添加核心功能和工具模块"

# 4. 改了一堆文件后统一提交
git add .
git commit -m "fix: 修复了三处边界条件处理问题"

# 5. 查看历史
git log --oneline
\end{lstlisting}

\subsection{查看历史 —— git log}

\begin{lstlisting}
# 基本查看
git log

# 最常用的精简模式
git log --oneline
# 输出：
# 3a4b5c6 修复了登录 bug
# 2b3c4d5 添加了用户注册功能
# 1a2b3c4 初始化项目

# 图形化查看所有分支
git log --graph --oneline --all
# 输出：
# * 3a4b5c6 (HEAD -> main) 修复了登录 bug
# | * 4b5c6d7 (feature) 添加了新页面
# |/
# * 2b3c4d5 添加了用户注册功能

# 查看每个提交的详细差异
git log -p

# 查看文件变更统计
git log --stat
# 输出：xxx.py | 10 +++++-----

# 按条件筛选
git log --author="张三"
git log --since="2024-01-01" --until="2024-12-31"
git log --grep="bug"
git log -3                     # 只看最近 3 次提交

# 查看某次提交的详细内容
git show 3a4b5c6               # 用提交哈希的前几位即可
\end{lstlisting}

\subsection{比较差异 —— git diff}

\texttt{git diff} 用来查看文件的具体修改内容，是 code review 和检查错误的好帮手。

\begin{lstlisting}
# 工作区 vs 暂存区（还没 add 的修改）
git diff

# 暂存区 vs 最近一次提交（准备要提交的内容）
git diff --staged              # 或 git diff --cached

# 工作区 vs 最近一次提交（所有未提交的修改）
git diff HEAD

# 两个分支之间的差异
git diff main feature

# 两次提交之间的差异
git diff 3a4b5c6 2b3c4d5

# 只看某个文件的差异
git diff -- README.md

# 忽略空格变化的差异
git diff -w
\end{lstlisting}

\textbf{输出解读}：
\begin{lstlisting}[numbers=none]
diff --git a/index.html b/index.html
index e69de29..40559da 100644
--- a/index.html              # a 版本
+++ b/index.html              # b 版本
@@ -0,0 +1,3 @@               # 从第0行开始，加了3行
+<html>                       # "+" 表示新增的行
+<head><title>Hello</title></head>
+</html>
\end{lstlisting}

\subsection{删除与移动}

\begin{lstlisting}
# 删除文件（同时在文件系统删除 + Git 中记录删除）
git rm app.log

# 仅从 Git 中移除跟踪，保留本地文件（常见于 .gitignore 问题修复）
git rm --cached config.yml

# 重命名或移动文件
git mv oldname.txt newname.txt
# 等价于：
# mv oldname.txt newname.txt
# git rm oldname.txt
# git add newname.txt
\end{lstlisting}

\subsection{撤销操作}

撤销是 Git 初学者最常疑问的地方，关键在于区分\textbf{撤销了什么区域}。

\subsubsection{场景一：文件改了但还没 add（工作区级撤销）}

\begin{lstlisting}
# 用仓库里的版本覆盖工作区的修改（危险！修改会丢失）
git restore README.md
# 旧式写法（功能一样）：
git checkout -- README.md
\end{lstlisting}

\subsubsection{场景二：已经 add 了但还没 commit（暂存区级撤销）}

\begin{lstlisting}
# 从暂存区撤出，但保留工作区的修改
git restore --staged README.md
# 旧式写法：
git reset HEAD README.md
\end{lstlisting}

\subsubsection{场景三：已经 commit 了但还没 push（仓库级撤销）}

\begin{lstlisting}
# --soft：只移动 HEAD，工作区和暂存区都不动
# 场景：提交后发现忘了加某个文件
git reset --soft HEAD~1
# 现在文件还在暂存区，你只需要重新 commit

# --mixed（默认）：移动 HEAD，清空暂存区但保留工作区
# 场景：想把提交拆成多个更小的提交
git reset --mixed HEAD~1

# --hard：完全撤销，丢弃所有修改（非常危险！）
git reset --hard HEAD~1
# 场景：提交了错误的代码，想完全不要这次修改
\end{lstlisting}

\texttt{HEAD~1} 表示 HEAD 的上一个提交，\texttt{HEAD~3} 表示往上数 3 个。

\subsubsection{场景四：已经 push 了（安全的方式）}

\begin{lstlisting}
# 不要用 reset！因为 reset 会改写历史，push 会被拒绝
# 正确做法：用 revert 创建一个"反提交"
git revert HEAD              # 撤销最近一次提交
git revert 3a4b5c6           # 撤销指定提交

# revert 会创建一个新的提交，内容是"把指定提交的修改反向应用"
# 这是完全安全的，因为历史没有被改写
\end{lstlisting}

\subsubsection{撤销操作选择速查表}

\begin{longtable}{|p{3.5cm}|p{4.5cm}|p{4cm}|}
\hline
\textbf{状态} & \textbf{推荐命令} & \textbf{安全性} \\
\hline
改了还没 add & \texttt{git restore <file>} & 危险，修改丢失 \\
\hline
add 了还没 commit & \texttt{git restore --staged <file>} & 安全 \\
\hline
commit 了还没 push & \texttt{git reset HEAD\textasciitilde1} & 安全（--hard 除外） \\
\hline
已经 push 了 & \texttt{git revert <commit>} & 最安全 \\
\hline
\end{longtable}

%====================================================================
\section{分支管理 —— Git 最强大的特性}
%====================================================================

\subsection{什么是分支？}

\textbf{分支本质上就是一个指向某次提交的指针}。

类比：想象你正在写一本书（main 分支），突然想尝试另一个写法，但不想影响原来的版本。
于是你复印了一份（创建分支），在复印件上修改。如果新写法好，就合并回原稿；
如果不好，扔掉复印件即可。

\subsection{为什么 Git 的分支那么轻量？}

\begin{itemize}[itemsep=3pt]
    \item SVN 创建分支 = 复制整个项目目录到另一个位置（慢，占空间）
    \item Git 创建分支 = 创建一个 41 字节的文件，里面只记录一个哈希值（瞬间完成）
    \item 所以 Git 鼓励\textbf{频繁创建和合并分支}
\end{itemize}

\texttt{HEAD} 是一个特殊指针，它指向\textbf{你当前在哪个分支上}。
当你切换分支时，HEAD 就指向另一个分支。

\begin{lstlisting}[numbers=none]
# 当前在 main 分支
HEAD -> main -> commit (abc123)

# 创建 feature 分支后
HEAD -> main -> commit (abc123)
        feature -> commit (abc123)   # 两个指针指向同一个提交

# 在 feature 上提交后
HEAD -> main -> commit (abc123)
        feature -> commit (def456)   # feature 前进了
\end{lstlisting}

\subsection{分支操作}

\begin{lstlisting}
# 查看分支
git branch                     # 本地分支，* 表示当前所在分支
git branch -r                  # 远程分支（如 origin/main）
git branch -a                  # 全部

# 创建分支
git branch feature-login       # 创建分支（基于当前所在位置）

# 切换分支
git switch feature-login       # 新版命令（推荐）
git checkout feature-login     # 传统命令

# 创建并切换（一行搞定）
git switch -c feature-login
git checkout -b feature-login  # 传统写法

# 重命名分支
git branch -m old-name new-name

# 删除分支
git branch -d feature-login    # 只能删除已合并的
git branch -D feature-login    # 强制删除（不管是否合并）

# 查看每个分支的最后一次提交
git branch -v

# 查看已合并到当前分支的分支
git branch --merged

# 查看未合并的分支
git branch --no-merged
\end{lstlisting}

\subsection{合并 —— git merge}

合并是把两个分支的工作成果整合在一起。

\subsubsection{Fast-Forward 合并（直线合并）}

场景：main 分支在创建 feature 后没有新的提交。

\begin{lstlisting}[numbers=none]
合并前:
  main:    A --- B
                \
  feature:        C --- D

合并（git merge feature）后:
  main:    A --- B --- C --- D
  # Git 只是把 main 指针往前移动了，没有创建新的合并提交
\end{lstlisting}

\begin{lstlisting}
# Fast-Forward 合并
git switch main
git merge feature              # main 直接移到 feature 的位置
\end{lstlisting}

\subsubsection{Three-Way 合并（分叉合并）}

场景：main 和 feature 都有新的提交。

\begin{lstlisting}[numbers=none]
合并前:
  main:    A --- B --- C
                \
  feature:        D --- E

合并后:
  main:    A --- B --- C --- F (merge commit)
                \         /
  feature:        D --- E
  # Git 会创建新的合并提交 F，它有两个父提交
\end{lstlisting}

\begin{lstlisting}
# 如果想禁用 Fast-Forward（显式创建合并提交）
git merge --no-ff feature      # 即使可以 Fast-Forward 也创建合并节点
# 在团队协作中 --no-ff 更常用，因为保留了分支历史
\end{lstlisting}

\subsection{合并冲突}

当两个分支修改了\textbf{同一个文件的同一行}时，Git 无法自动判断哪个是正确版本，
需要你手动解决。

\subsubsection{制造冲突的示例}

\begin{lstlisting}[numbers=none]
1. 在 main 分支上修改 README.md 的第一行：
   echo "Main Branch Line" > README.md
   git add . && git commit -m "main 修改"

2. 在 feature 分支上修改同一个文件的同一行：
   git switch -c feature
   echo "Feature Branch Line" > README.md
   git add . && git commit -m "feature 修改"

3. 切换回 main 合并 feature：
   git switch main
   git merge feature
   # 输出：CONFLICT (content): Merge conflict in README.md
\end{lstlisting}

\subsubsection{冲突标记}

\begin{lstlisting}
# 冲突文件里的内容：
<<<<<<< HEAD
Main Branch Line              # 当前分支（main）的内容
=======
Feature Branch Line           # 要合并分支（feature）的内容
>>>>>>> feature

# 你需要手动编辑文件，决定保留什么
\end{lstlisting}

\subsubsection{解决冲突的完整步骤}

\begin{lstlisting}
# 1. 查看哪些文件冲突了
git status
# 输出：both modified:   README.md

# 2. 编辑冲突文件，删除冲突标记，保留你想要的内容
# 比如你决定两个都要：
# Main Branch Line
# Feature Branch Line

# 3. 标记为已解决
git add README.md

# 4. 完成合并
git commit
# Git 会自动生成合并提交消息，你也可以修改

# 5. 如果发现冲突太难解决，想回到合并前
git merge --abort
\end{lstlisting}

\subsection{Rebase（变基）}

Rebase 的原理是：把当前分支的提交\textbf{移植}到另一个分支的顶端。

\subsubsection{Rebase vs Merge}

\begin{lstlisting}[numbers=none]
# Merge 的方式（保留分叉历史）：
# main:   A --- B --- C --- F (merge commit)
#               \         /
# feature:       D --- E

# Rebase 的方式（线性历史）：
# main:   A --- B --- C
#                         \
# feature:                 D' --- E'  (新提交，哈希变了)
# 然后把 main fast-forward 到 E'
# main:   A --- B --- C --- D' --- E'
\end{lstlisting}

\begin{lstlisting}
# Rebase 操作
git switch feature
git rebase main                # 把 feature 的提交搬到 main 的最新提交之后

# 如果 rebase 过程中有冲突：
# 1. 解决冲突
# 2. git add <file>
# 3. git rebase --continue
# 或者 git rebase --abort 放弃 rebase
\end{lstlisting}

\subsubsection{交互式 Rebase —— 改写历史的瑞士军刀}

\begin{lstlisting}
# 修改最近 3 次提交
git rebase -i HEAD~3
# 这会打开一个编辑器，显示：
pick 1a2b3c4 第一个提交
pick 2b3c4d5 第二个提交
pick 3c4d5e6 第三个提交

# 把 pick 改成对应的命令：
\end{lstlisting}

\textbf{可用命令}：
\begin{itemize}[itemsep=3pt]
    \item \texttt{pick}（或 p）— 保留该提交
    \item \texttt{reword}（或 r）— 只改提交消息，不改内容
    \item \texttt{edit}（或 e）— 修改提交内容
    \item \texttt{squash}（或 s）— 合并到上一个提交（保留提交消息）
    \item \texttt{fixup}（或 f）— 合并到上一个提交（丢弃消息）
    \item \texttt{drop}（或 d）— 删除该提交
    \item \texttt{exec}（或 x）— 在此处执行 shell 命令
\end{itemize}

\textbf{常见用法}：把一堆零碎的提交合并成一个整洁的提交。

\begin{lstlisting}[numbers=none]
# 改前（杂乱的提交历史）：
fix typo
add comment
finish feature

# rebase -i 改成：
pick finish feature
fixup fix typo                  # 合并到上一个
fixup add comment               # 合并到上一个

# 改后（整洁的历史）：
finish feature
\end{lstlisting}

\subsubsection{黄金法则 —— 永远记住！}

\begin{center}
\colorbox{yellow}{\parbox{0.9\textwidth}{\centering
\textbf{绝不}对\textbf{已经推送到远程仓库}的提交执行 rebase！
}}
\end{center}

因为 rebase 会改变提交的哈希值，如果你 rebase 了已推送的提交，
下次 \texttt{git push} 会被拒绝（因为历史不一致）。
即使强制推送（\texttt{git push --force}）成功，队友的仓库也会陷入混乱。

\subsection{Cherry-Pick —— 精准挑选提交}

有时候你不想合并整个分支，只想\textbf{挑某个分支中的一个提交}拿过来。

\begin{lstlisting}
# 把 feature 分支上的某个提交（哈希 abc123）应用到当前分支
git cherry-pick abc123

# 一次挑多个
git cherry-pick abc123 def456

# 如果有冲突，解决后：
git cherry-pick --continue
# 或者放弃：
git cherry-pick --abort
\end{lstlisting}

使用场景：在 main 分支上发现了 feature 分支中的一个 bug 修复提交，
没必要合并整个 feature 分支，用 cherry-pick 把那个修复提交拿过来即可。

%====================================================================
\section{远程仓库 —— 和世界协作}
%====================================================================

\subsection{远程仓库的概念}

远程仓库（remote）是存放在\textbf{服务器}上的 Git 仓库。常见的有：

\begin{itemize}[itemsep=3pt]
    \item GitHub（全球最大的代码托管平台）
    \item GitLab（企业级自托管方案）
    \item Gitee（国内代码托管平台）
\end{itemize}

远程仓库让你的代码：
\begin{enumerate}[itemsep=3pt]
    \item 有了\textbf{备份}——电脑坏了也不怕
    \item 可以\textbf{协作}——团队多人共同开发
    \item 方便\textbf{分享}——开源项目
\end{enumerate}

\subsection{克隆已有仓库}

\begin{lstlisting}
# 克隆到当前目录（会创建同名文件夹）
git clone https://github.com/username/repo.git

# 克隆到指定目录名
git clone https://github.com/username/repo.git my-folder

# 克隆指定分支
git clone -b feature https://github.com/username/repo.git
\end{lstlisting}

\subsection{管理远程仓库}

\begin{lstlisting}
# 查看已经配置的远程仓库
git remote -v
# 输出：
# origin  https://github.com/username/repo.git (fetch)
# origin  https://github.com/username/repo.git (push)

# 添加新的远程仓库
git remote add origin https://github.com/username/repo.git

# 删除远程仓库
git remote remove origin

# 重命名远程仓库
git remote rename origin upstream

# 查看远程仓库的详细信息
git remote show origin
\end{lstlisting}

\texttt{origin} 是远程仓库的默认名字，就像 \texttt{main} 是默认分支名一样。
你可以改，但约定俗成就用 origin。

\subsection{推送 —— git push}

\begin{lstlisting}
# 首次推送（必须用 -u 建立关联）
git push -u origin main
# -u 是 --set-upstream 的缩写
# 之后就可以直接用 git push 了

# 推送其他分支
git push origin feature-login

# 推送所有分支
git push --all

# 删除远程分支
git push origin --delete feature-login

# 强制推送（非常危险！会覆盖远程历史）
git push --force origin main
# 只在确定没有其他人使用这个分支时才用
\end{lstlisting}

\subsection{拉取 —— git pull 与 git fetch}

这是初学者最容易混淆的地方。

\begin{itemize}[itemsep=5pt]
    \item \textbf{git fetch} — \textcolor{blue}{只下载，不合并}
    \begin{itemize}
        \item 从远程下载最新的提交到你的本地仓库
        \item 但\textbf{不会修改你的工作区}
        \item 你需要手动用 \texttt{git merge origin/main} 来合并
        \item 安全，可以随时查看远程有什么变化后再决定
    \end{itemize}

    \item \textbf{git pull} — \textcolor{blue}{下载并自动合并}
    \begin{itemize}
        \item = \texttt{git fetch} + \texttt{git merge}
        \item 一步到位
        \item 但可能会产生多余的合并提交
    \end{itemize}

    \item \textbf{git pull --rebase} — \textcolor{blue}{推荐！}
    \begin{itemize}
        \item = \texttt{git fetch} + \texttt{git rebase}
        \item 拉取远程更新后，把你本地的提交放在最上面
        \item 保持历史线性整洁
    \end{itemize}
\end{itemize}

\begin{lstlisting}
# 查看远程有更新但不去合并
git fetch origin

# 查看本地和远程的区别
git log main..origin/main --oneline

# 确认无误后合并
git merge origin/main

# 或者一步到位（推荐用 rebase）
git pull --rebase origin main

# 如果 pull 有冲突，按常规冲突解决流程走
\end{lstlisting}

\subsection{标准协作流程}

\begin{lstlisting}
# 假设你在一个团队项目中

# 1. 第一天：克隆项目
git clone https://github.com/team/project.git
cd project

# 2. 创建功能分支开发
git switch -c feature-user-page

# 3. 开发过程中，每天上班先同步 main
git switch main
git pull --rebase               # 获取最新的 main
git switch feature-user-page
git rebase main                 # 你的分支也同步 main 的更新

# 4. 提交你的工作
git add .
git commit -m "feat: 添加用户页面布局"
git push origin feature-user-page

# 5. 去 GitHub 上创建 Pull Request（PR）
# 等待 code review...

# 6. PR 合并后，删除本地分支
git switch main
git pull --rebase
git branch -d feature-user-page
\end{lstlisting}

\subsection{SSH 与 HTTPS 的区别}

\begin{longtable}{|p{3.5cm}|p{5cm}|p{5cm}|}
\hline
\textbf{特性} & \textbf{HTTPS} & \textbf{SSH} \\
\hline
克隆地址 & \texttt{https://github.com/...} & \texttt{git@github.com:...} \\
\hline
认证方式 & 每次 push 需要输入用户名密码 & 密钥对认证 \\
\hline
配置难度 & 简单，无需额外配置 & 需要生成 SSH 密钥并添加到 GitHub \\
\hline
推荐场景 & 快速开始、偶尔使用 & 日常频繁使用 \\
\hline
\end{longtable}

\begin{lstlisting}
# 配置 SSH（以 GitHub 为例）
# 1. 生成密钥对
ssh-keygen -t ed25519 -C "your_email@example.com"

# 2. 查看公钥
cat ~/.ssh/id_ed25519.pub

# 3. 复制公钥，添加到 GitHub Settings > SSH and GPG keys

# 4. 测试连接
ssh -T git@github.com
# 输出：Hi username! You've successfully authenticated...
\end{lstlisting}

%====================================================================
\section{标签管理 —— 标记重要版本}
%====================================================================

标签用来给某个特定的提交打上\textbf{有意义的名称}，通常用于版本发布。

\subsection{轻量标签 vs 附注标签}

\begin{itemize}[itemsep=5pt]
    \item \textbf{轻量标签（lightweight）}
    \begin{itemize}
        \item 只是一个指向提交的指针，像一个不会移动的分支
        \item 创建：\texttt{git tag v1.0.0}
        \item 适合临时标记或私用
    \end{itemize}

    \item \textbf{附注标签（annotated）}
    \begin{itemize}
        \item 作为完整对象存储在 Git 数据库中
        \item 包含：创建者、邮箱、日期、标签消息
        \item 可以用 GPG 签名验证真伪
        \item 创建：\texttt{git tag -a v1.0.0 -m "发布 1.0.0"}
        \item \textbf{推荐用于正式发布}
    \end{itemize}
\end{itemize}

\begin{lstlisting}
# 创建附注标签（推荐）
git tag -a v1.0.0 -m "正式发布 1.0.0"

# 为历史提交打标签
git tag -a v0.9.0 abc123 -m "发布前的 beta 版本"

# 查看标签
git tag                           # 列出所有标签
git tag -l "v1.*"                 # 搜索 v1.x 的标签
git show v1.0.0                   # 查看标签详细信息

# 推送标签
git push origin v1.0.0            # 推送一个标签
git push origin --tags            # 推送所有未推送的标签

# 删除标签
git tag -d v1.0.0                 # 删除本地
git push origin --delete v1.0.0   # 删除远程

# 切换到标签时创建分支（因为标签是只读的）
git switch -c release-v1 v1.0.0
\end{lstlisting}

\textbf{语义化版本控制（SemVer）}：
\begin{lstlisting}[numbers=none]
vMAJOR.MINOR.PATCH
   |     |     |
   |     |     +-- PATCH: 修复 bug，向后兼容
   |     +-- MINOR: 添加功能，向后兼容
   +-- MAJOR: 不兼容的重大变更
\end{lstlisting}

%====================================================================
\section{储藏（Stash）—— 临时保存工作现场}
%====================================================================

\subsection{什么时候需要 Stash？}

你正在开发一个功能，改了 5 个文件，突然需要切换到另一个分支修一个紧急 bug。
但这些改动还没完成，不想提交。怎么办？用 \texttt{git stash}。

\begin{lstlisting}
# 开发中...突然有紧急任务
git stash                        # 保存工作区所有修改，工作区恢复干净
# 输出：Saved working directory and index state WIP on main: abc123 ...

# 切换到其他分支修 bug
git switch hotfix
# ...修完 bug 后...

# 切回来恢复工作现场
git switch main
git stash pop                    # 恢复并删除栈顶的储藏
\end{lstlisting}

\subsection{Stash 的更多用法}

\begin{lstlisting}
# 储藏时加说明（方便查找）
git stash save "用户列表页 WIP"

# 查看所有储藏
git stash list
# 输出：
# stash@{0}: On main: 用户列表页 WIP
# stash@{1}: On main: 登录功能开发中

# 恢复但不删除
git stash apply                  # 恢复最近一个
git stash apply stash@{1}        # 恢复指定的

# 删除指定储藏
git stash drop stash@{0}

# 清空所有
git stash clear

# 储藏未跟踪的文件（默认只会储藏已跟踪的文件）
git stash -u                     # -u = --include-untracked

# 从某个储藏中创建新分支（如果恢复时有冲突，推荐用这个）
git stash branch new-branch stash@{0}
\end{lstlisting}

Stash 是一个\textbf{栈结构}，后保存的在栈顶（stash@\{0\}）。

%====================================================================
\section{.gitignore —— 告诉 Git 忽略什么}
%====================================================================

\subsection{为什么要忽略文件？}

你的项目里有很多文件是不需要被版本控制的：
\begin{itemize}[itemsep=3pt]
    \item 编译产物（\texttt{.pyc}、\texttt{.o}、\texttt{.exe}）
    \item 依赖目录（\texttt{node\_modules/}、\texttt{vendor/}）
    \item 系统文件（\texttt{.DS\_Store}、\texttt{Thumbs.db}）
    \item 环境配置（\texttt{.env}、\texttt{*.local}）
    \item IDE 配置（\texttt{.vscode/}、\texttt{.idea/}）
    \item 日志文件（\texttt{*.log}）
\end{itemize}

\textbf{不忽略这些文件会导致}：
\begin{itemize}[itemsep=3pt]
    \item 仓库臃肿（\texttt{node\_modules} 可能有几百 MB）
    \item 提交混乱（编译文件每个版本都变化）
    \item 安全风险（\texttt{.env} 可能包含密码）
\end{itemize}

\subsection{语法规则}

\begin{lstlisting}
# 注释以 # 开头

# 1. 忽略特定文件
.DS_Store              # 忽略当前目录下的 .DS_Store
.env                   # 忽略 .env

# 2. 通配符
*.log                  # 忽略所有 .log 结尾的文件
debug.log              # 但 .gitignore 也可以针对具体文件

# 3. 忽略目录（注意结尾的 /）
node_modules/          # 忽略 node_modules 目录及其所有内容
build/                 # 忽略 build 目录

# 4. 取反（不忽略某个特定文件）
*.log                  # 忽略所有 .log
!important.log         # 但不忽略 important.log

# 5. 只忽略根目录下的（不匹配子目录）
/ToDo.txt              # 只忽略根目录下的 ToDo.txt
                       # 不忽略 src/ToDo.txt

# 6. 忽略所有目录下某类文件，但不包括子目录中的
*/build/               # 匹配 a/build/ 但不匹配 a/b/build/

# 7. 忽略除某文件外的所有文件（白名单模式）
/*
!src/
!src/**/*.py           # 只跟踪 src 目录下的 .py 文件
\end{lstlisting}

\subsection{常见语言/项目的 .gitignore 模板}

\begin{lstlisting}
# Python 项目
__pycache__/
*.py[cod]
*.so
*.egg-info/
dist/
build/
.venv/
env/

# Node.js 项目
node_modules/
npm-debug.log*
dist/
build/
.next/

# Java 项目
*.class
*.jar
*.war
target/
!.mvn/wrapper/maven-wrapper.jar

# LaTeX 项目（你的笔记项目用这个）
*.aux
*.log
*.out
*.toc
*.fdb_latexmk
*.fls
*.synctex.gz
\end{lstlisting}

\textbf{重要提醒}：如果文件已经被 Git 跟踪了，再加入 .gitignore 是不生效的！
需要用 \texttt{git rm --cached} 先解除跟踪。

%====================================================================
\section{实用技巧与最佳实践}
%====================================================================

\subsection{别名（Alias）—— 少打字}

\begin{lstlisting}
git config --global alias.st status
git config --global alias.co checkout
git config --global alias.br branch
git config --global alias.ci commit
git config --global alias.unstage "reset HEAD"
git config --global alias.last "log -1 HEAD"

# 设置完后，git st 就等于 git status
git st

# 高级别名：显示漂亮的日志
git config --global alias.tree "log --graph --oneline --all"
# 用法：git tree

# 更美观的版本
git config --global alias.lg "log --graph --pretty=format:'%Cred%h%Creset -%C(yellow)%d%Creset %s %Cgreen(%cr) %C(bold blue)<%an>%Creset' --abbrev-commit"
\end{lstlisting}

\subsection{查看日志的美化方式}

\begin{lstlisting}
# 一行一个提交，带图形
git log --graph --oneline --all

# 自定义格式（可以直接设为别名）
git log --graph --pretty=format:'%Cred%h%Creset -%C(yellow)%d%Creset %s %Cgreen(%cr) %C(bold blue)<%an>%Creset' --abbrev-commit
# 效果：红色哈希 - 黄色分支名 提交消息 绿色时间 <蓝色作者>
\end{lstlisting}

\subsection{常见场景速查}

\begin{lstlisting}
# 1. 想在 main 上修改，但已经开发了一半
git stash                      # 保存当前修改
git switch main                # 切回 main
# ...做完紧急修改并提交后...
git switch feature             # 切回功能分支
git stash pop                  # 恢复工作现场

# 2. 提交到了 main，但应该提交到 feature 分支
git branch feature             # 在当前提交创建 feature 分支
git reset --hard HEAD~1        # main 退回一步
git switch feature             # 切到 feature（提交还在）

# 3. 想修改倒数第二次的提交消息
git rebase -i HEAD~3
# 把第二行的 pick 改成 r（reword），保存退出后修改消息

# 4. 把多个零碎提交合并成一个
git rebase -i HEAD~5
# 除了第一个 pick，后面的都改成 s（squash）或 f（fixup）

# 5. 想看看某个文件的历史修改
git log -p -- README.md        # 看 README.md 的完整修改历史
git blame README.md            # 每行代码是谁改的

# 6. 批量修改最近的提交信息
git rebase -i HEAD~5
# 把要改的 pick 改成 r（reword）

# 7. 找回误删的分支
git reflog                     # 查看所有 HEAD 移动的历史
# 输出：abc123 HEAD@{0}: Branch: deleted
git switch -c recovered-branch abc123  # 在哈希处创建新分支
\end{lstlisting}

\subsection{救命的 reflog}

\texttt{git reflog} 是 Git 的\textbf{操作日志}，记录了 HEAD 的所有移动记录。
即使你 reset、rebase 搞乱了，也能通过 reflog 找回。

\begin{lstlisting}
# 查看最近的操作记录
git reflog
# 输出：
# abc123 (HEAD -> main) HEAD@{0}: commit: 修复登录 bug
# def456 HEAD@{1}: reset: moving to HEAD~1
# ghi789 HEAD@{2}: commit: 测试功能
# jkl012 HEAD@{3}: commit: 添加用户页面
# 注意：reflog 只记录本地操作，每个仓库的 reflog 都不同

# 最实用的场景：找回 reset --hard 丢失的提交
git reset --hard HEAD~2        # 完了，回退太远了！
git reflog                     # 查看历史
git reset --hard ghi789        # 恢复回去
\end{lstlisting}

\subsection{二分查找 —— git bisect}

当你发现一个 bug，但不知道是哪个提交引入的，可以用 bisect 二分查找。

\begin{lstlisting}
# 假设你有 100 个提交，手动排查要查 100 次
# 用二分法最多查 log2(100) = 7 次

# 开始二分查找
git bisect start

# 标记当前版本有 bug
git bisect bad

# 标记某个历史版本没有 bug（比如 50 个提交前的版本）
git bisect good abc123

# Git 会自动切换到中间的某个提交
# 你测试这个版本，然后标记：
git bisect good                # 如果这个版本没有问题
git bisect bad                 # 如果这个版本有 bug

# Git 会继续二分切换到下一个要测试的提交
# 重复 good/bad...

# 最终 Git 会告诉你是哪个提交引入的 bug
# abcdef is the first bad commit

# 结束 bisect
git bisect reset
\end{lstlisting}

\textbf{自动化 bisect}：如果项目有测试脚本，可以让 Git 自动运行：

\begin{lstlisting}
git bisect start
git bisect bad
git bisect good abc123
git bisect run python test.py   # 自动：测试通过 = good，失败 = bad
\end{lstlisting}

\subsection{子模块（Submodule）}

当一个 Git 仓库需要引用另一个 Git 仓库时（比如项目引用了某个第三方库），
可以使用 submodule。

\begin{lstlisting}
# 添加子模块
git submodule add https://github.com/lib/lib.git libs/lib

# 克隆包含子模块的项目
git clone --recurse-submodules https://github.com/project.git

# 更新子模块
git submodule update --remote

# 注意：子模块默认不会自动更新，需要手动更新
\end{lstlisting}

\subsection{提交信息规范 —— Conventional Commits}

\textbf{为什么要规范提交消息？}
\begin{itemize}[itemsep=3pt]
    \item 方便生成 changelog
    \item 更好的搜索和过滤
    \item 自动化的版本号管理
\end{itemize}

\textbf{标准格式}：

\begin{lstlisting}
<type>(<scope>): <subject>

<body>

<footer>
\end{lstlisting}

\textbf{常用类型}：
\begin{longtable}{|p{3cm}|p{4cm}|p{5cm}|}
\hline
\textbf{类型} & \textbf{含义} & \textbf{示例} \\
\hline
feat & 新功能 & \texttt{feat: 添加用户注册功能} \\
\hline
fix & 修复 bug & \texttt{fix: 修复密码验证越界问题} \\
\hline
docs & 文档变更 & \texttt{docs: 更新 API 文档} \\
\hline
style & 代码格式 & \texttt{style: 格式化代码缩进} \\
\hline
refactor & 重构 & \texttt{refactor: 拆分用户验证模块} \\
\hline
test & 测试相关 & \texttt{test: 添加登录测试用例} \\
\hline
chore & 构建/工具 & \texttt{chore: 更新依赖版本} \\
\hline
perf & 性能优化 & \texttt{perf: 优化图片加载速度} \\
\hline
\end{longtable}

\textbf{好例子 vs 坏例子}：

\begin{lstlisting}[numbers=none]
# 坏例子
fix bug                             # 太模糊
update file                         # 说了等于没说
asdfgh                              # 毫无意义

# 好例子
fix(login): 修复密码为空时程序崩溃
feat(user): 添加用户头像上传功能
docs(readme): 更新安装步骤，添加 Windows 系统支持
refactor(db): 提取数据库连接为独立模块
chore: 更新 ESLint 配置，统一代码风格
\end{lstlisting}

\subsection{Git 钩子（Hooks）}

Git 钩子是在特定事件发生时自动触发的脚本，放在 \texttt{.git/hooks/} 目录下。

常见钩子：
\begin{itemize}[itemsep=3pt]
    \item \texttt{pre-commit}：提交前运行，用于检查代码风格、运行测试
    \item \texttt{commit-msg}：检查提交消息格式是否符合规范
    \item \texttt{pre-push}：推送前运行
    \item \texttt{post-merge}：合并后自动执行
\end{itemize}

\begin{lstlisting}
# 示例：pre-commit 钩子（.git/hooks/pre-commit）
#!/bin/sh
# 检查是否有未解决的冲突标记
if grep -r "^<<<<<<<\|^>>>>>>>" --include="*.py" .; then
    echo "错误：存在未解决的合并冲突！"
    exit 1
fi
\end{lstlisting}

注意：\texttt{.git/hooks/} 下的钩子\textbf{不会被提交}。团队共享钩子需要借助工具（如 Husky）。

\subsection{工作流模型}

\subsubsection{Git Flow —— 适合有固定发布周期的项目}

\begin{lstlisting}[numbers=none]
main ────────●──────────────●────────── (生产发布)
              \            /
develop ──────●──●──●────●──────────── (开发主线)
               \    /  \    /
feature/*       ●──●    ●──●           (功能开发)

hotfix/* ──────●──●                    (紧急修复，从 main 切出)
\end{lstlisting}

\begin{itemize}[itemsep=3pt]
    \item \texttt{main} — 生产分支，始终处于可部署状态
    \item \texttt{develop} — 开发主分支，集成所有功能
    \item \texttt{feature/*} — 从 develop 创建，完成后合并回 develop
    \item \texttt{release/*} — 发布前准备，只修 bug 不加功能
    \item \texttt{hotfix/*} — 从 main 创建，修复后合并回 main 和 develop
\end{itemize}

\textbf{优点}：流程清晰，适合大项目。\\
\textbf{缺点}：分支多且复杂，不适合持续部署。

\subsubsection{GitHub Flow —— 适合持续部署}

\begin{lstlisting}[numbers=none]
main ────●────────●───────────●────── (随时可部署)
          \        /           /
feature    ●──●──●            /
                \            /
fix              ●──●──●──●
\end{lstlisting}

\begin{itemize}[itemsep=3pt]
    \item \texttt{main} 是唯一长期分支，始终可部署
    \item 从 main 创建特性分支
    \item 提交修改，频繁推送
    \item 创建 Pull Request 进行 Code Review
    \item Review 通过后合并到 main
    \item 合并后立即部署
\end{itemize}

\textbf{优点}：极其简单，适合 CI/CD。\\
\textbf{缺点}：不适合需要多个发布版本同时维护的项目。

\subsubsection{GitLab Flow}

GitLab 的变体，在 GitHub Flow 基础上增加了\textbf{环境分支}：

\begin{lstlisting}[numbers=none]
main ────●────────●───────────────
           \        \
staging      ●──────●──────────────
                     \
production            ●──────────
\end{lstlisting}

环境分支（\texttt{staging}、\texttt{production}）确保代码经过充分测试再上线。

%====================================================================
\section{内部原理速览 —— Git 是如何工作的？}
%====================================================================

\subsection{.git 目录全解析}

\begin{lstlisting}[numbers=none]
.git/
|-- HEAD                  # 指向当前分支（如 ref: refs/heads/main）
|-- config                # 本仓库的 Git 配置
|-- description           # 仓库描述（给 GitWeb 用的）
|-- index                 # 暂存区（二进制文件）
|-- ORIG_HEAD             # 某些危险操作前 HEAD 的备份
|-- FETCH_HEAD            # 最近一次 fetch 的来源信息
|
|-- objects/              # 对象存储（Git 的核心！）
|   |-- 1a/               # 哈希前两位做目录名
|   |   `-- b2c3d4e...    # 文件名 = 后 38 位
|   |-- pack/             # 打包文件（压缩优化）
|   `-- info/             # 包文件索引
|
|-- refs/                 # 引用
|   |-- heads/            # 本地分支（每个文件存一个哈希）
|   |   |-- main
|   |   `-- feature
|   |-- tags/             # 标签
|   |   `-- v1.0.0
|   `-- remotes/          # 远程分支
|       `-- origin/
|           `-- main
|
`-- hooks/                # Git 钩子脚本示例
    |-- pre-commit.sample
    |-- commit-msg.sample
    `-- ...
\end{lstlisting}

\subsection{底层命令实验}

理解 Git 的最佳方式是用底层命令亲手操作：

\begin{lstlisting}
# 1. 手动创建一个 blob 对象
echo "Hello Git" | git hash-object --stdin -w
# 输出：哈希值（同时在 .git/objects/ 下创建了文件）

# 2. 查看对象类型
git cat-file -t <hash>
# 输出：blob

# 3. 查看对象内容
git cat-file -p <hash>
# 输出：Hello Git

# 4. 查看暂存区的内容
git ls-files --stage

# 5. 查看不同引用指向的提交
git rev-parse HEAD           # 查看 HEAD 的哈希
git rev-parse main           # 查看 main 分支的哈希
git rev-parse v1.0.0         # 查看标签的哈希
\end{lstlisting}

\subsection{理解 git add、git commit 的底层操作}

\begin{lstlisting}[numbers=none]
当你执行 git add README.md 时，Git 实际上做了：
1. 用 README.md 的内容创建一个 blob 对象
2. 更新 index（暂存区），记录 README.md -> 该 blob

当你执行 git commit 时，Git 实际上做了：
1. 基于 index 创建一个 tree 对象（记录目录结构）
2. 创建一个 commit 对象，指向该 tree，并记录父提交等元信息
3. 将 HEAD 指向这个新的 commit
\end{lstlisting}

%====================================================================
\section{常见问题 FAQ}
%====================================================================

\begin{enumerate}[itemsep=10pt]

    \item \textbf{Q: 不小心把密码提交到仓库了怎么办？}
    \begin{itemize}
        \item 如果还没有 push：用 \texttt{git reset} 撤销提交
        \item 如果已经 push：立即修改密码！然后 \texttt{git push --force} 覆盖历史（但 GitHub/GitLab 上可能已被爬虫抓取）
        \item 真正的解决方案：以后用 \texttt{.gitignore} + 环境变量
    \end{itemize}

    \item \textbf{Q: 为什么我 git push 被拒绝了？}
    \begin{itemize}
        \item 最常见原因：远程仓库有新的提交你没拉到
        \item 解决方案：\texttt{git pull --rebase} 后再 push
    \end{itemize}

    \item \textbf{Q: merge 和 rebase 到底选哪个？}
    \begin{itemize}
        \item 个人开发/小团队：用 rebase 保持历史整洁
        \item 团队协作/公共分支：用 merge --no-ff 保留分支信息
        \item 规则：对\textbf{尚未推送}的提交用 rebase，已推送的用 merge
    \end{itemize}

    \item \textbf{Q: git pull 时冲突了怎么办？}
    \begin{itemize}
        \item 和普通冲突一样解决，解决后 \texttt{git add} + \texttt{git commit}
        \item 如果不想解决：\texttt{git merge --abort} 回到 pull 前的状态
    \end{itemize}

    \item \textbf{Q: 想放弃当前所有修改怎么办？}
    \begin{itemize}
        \item \texttt{git restore .} 撤销所有未 add 的修改
        \item \texttt{git reset --hard HEAD} 撤销所有未提交的修改
        \item \texttt{git clean -fd} 删除所有未跟踪的文件和目录
        \item 组合使用：\texttt{git reset --hard HEAD \&\& git clean -fd}
        \item \textbf{以上操作均不可恢复，谨慎使用！}
    \end{itemize}

    \item \textbf{Q: .gitignore 为什么不生效？}
    \begin{itemize}
        \item 检查文件名是否拼写正确（不是 .gitingore）
        \item 检查是否已经被 Git 跟踪了（用 \texttt{git ls-files} 查看）
        \item 如果已跟踪：先用 \texttt{git rm --cached} 解除跟踪
    \end{itemize}

\end{enumerate}

%====================================================================
\section{总结}
%====================================================================

Git 学习核心要点：

\begin{enumerate}[itemsep=8pt]
    \item \textbf{理解三个区域}：工作区、暂存区、仓库是 Git 的核心抽象。
    所有操作都是这三个区域之间的数据移动。

    \item \textbf{分支是指针}：创建和切换分支代价极低，所以\textbf{大胆创建分支}，
    这是 Git 最大的优势。

    \item \textbf{小步提交，频繁提交}：每次提交只做一件事，提交消息写清楚
    \texttt{为什么}。这让你可以精准地回退到任何状态。

    \item \textbf{拉取用 rebase，合并用 merge}：
    \texttt{pull --rebase} 保持本地历史整洁，\texttt{merge --no-ff} 保留分支协作信息。

    \item \textbf{推送前检查}：
    \texttt{git status} 看状态，\texttt{git diff} 看内容，确认无误再推送。

    \item \textbf{.gitignore 尽早配置}：项目一开始就配好，避免跟踪不必要的文件。

    \item \textbf{遇到问题先 git status}：这是诊断一切问题的起点。

    \item \textbf{不要惊慌}：Git 的大部分操作都可以撤销。
    \texttt{reflog} 是你的最后一道防线，几乎可以找回一切。

    \item \textbf{多用命令行}：GUI 工具虽然方便，但理解 Git 最好的方式还是命令行。

    \item \textbf{练习是最好的老师}：开一个新仓库，故意制造冲突、练习 rebase、
    试试 reset 和 revert。用多了自然熟练。
\end{enumerate}

\begin{center}
\vspace{1cm}
\textcolor{gray}{\small 学习 Git 最好的方式就是多用！祝编码愉快！}
\end{center}

\end{document}
